\documentclass[11pt,a4paper]{article}
% Compile with XeLaTeX. All nonstandard text fonts are bundled in fonts/.
% The layout follows paper/main.tex; Russian text requires Unicode fonts.
\usepackage{fontspec}
\usepackage{polyglossia}
\setmainlanguage{russian}
\setotherlanguage{english}
\setmainfont{NimbusRoman}[Path=fonts/,Extension=.otf,
 UprightFont=*-Regular,BoldFont=*-Bold,ItalicFont=*-Italic,
 BoldItalicFont=*-BoldItalic]
\setmonofont{NimbusMonoPS}[Path=fonts/,Extension=.otf,
 UprightFont=*-Regular,BoldFont=*-Bold,Scale=0.85]
\usepackage{amsmath,amsthm,mathtools}
\usepackage{unicode-math}
\setmathfont{Latin Modern Math}
\usepackage[margin=24mm,headheight=15pt]{geometry}
\usepackage{microtype}
\usepackage{xcolor,booktabs,longtable,array,enumitem,fancyhdr,needspace}
\usepackage{xurl}
\usepackage[unicode,colorlinks=true,linkcolor=blue!45!black,
 citecolor=blue!45!black,urlcolor=blue!45!black]{hyperref}
\hypersetup{pdftitle={Appendix 2: Cremona groups and polynomial rectification},
 pdfauthor={Ivan Dorofeev and Alexander Frolov},pdfsubject={Research branch under verification; 28 October 2026}}
\pagestyle{fancy}\fancyhf{}
\fancyhead[L]{\small Appendix 2: группы Кремоны}
\fancyhead[R]{\small На стадии проверки}
\fancyfoot[C]{\thepage}
\renewcommand{\headrulewidth}{0.3pt}
\setlength{\parindent}{0pt}\setlength{\parskip}{5pt}
\setlength{\emergencystretch}{3em}
\setlist{nosep,leftmargin=*}
\newtheorem{theorem}{Теорема}[section]
\newtheorem{proposition}[theorem]{Предложение}
\newtheorem{lemma}[theorem]{Лемма}
\newtheorem{corollary}[theorem]{Следствие}
\theoremstyle{definition}
\newtheorem{definition}[theorem]{Определение}
\newtheorem{question}[theorem]{Вопрос}
\theoremstyle{remark}
\newtheorem{remark}[theorem]{Замечание}
\newcommand{\A}{{\mathbb A}}
\newcommand{\Q}{{\mathbb Q}}
\newcommand{\C}{{\mathbb C}}
\newcommand{\Z}{{\mathbb Z}}
\newcommand{\R}{{\mathbb R}}
\newcommand{\Aut}{\operatorname{Aut}}
\newcommand{\Cr}{\operatorname{Cr}}
\newcommand{\Tame}{\operatorname{Tame}}
\newcommand{\id}{\mathrm{id}}
\newcommand{\Stab}{\operatorname{Stab}}
\newcommand{\Frac}{\operatorname{Frac}}
\title{\vspace{-12mm}\LARGE\bfseries Appendix 2\\[4pt]
 Группы Кремоны и полиномиальное выпрямление\\
 вложенных прямых: результаты и программа исследования}
\author{Ivan Dorofeev and Alexander Frolov\\[3pt]\normalsize Researchers\\
 \normalsize\href{mailto:qft98romanov@gmail.com}{\texttt{qft98romanov@gmail.com}}}
\date{28 октября 2026 года}

\begin{document}
\maketitle
\begin{abstract}
Исследовательское продолжение явных выпрямлений трилистника Шастри и
гармонической восьмёрки над $\Q(\sqrt{-3})$. Задача рассматривается через
различие между бирациональной и полиномиальной эквивалентностью, относительные
автоморфизмы и устранение полюсов. Приведены рабочие доказательства дикости
выбранных выпрямителей, их принадлежности стандартной кремоновой подгруппе
и стабильной ручности. Построено семейство переносов и дан точный учёт
того, какие конкретные вопросы получают ответ, какие результаты заимствованы
из литературы и что остаётся открытым. Символьные сертификаты используют
точную арифметику и сопровождаются воспроизводимым кодом.
\end{abstract}

\paragraph{Благодарности.}
Авторы выражают искреннюю благодарность профессору Алексею Канель-Белову
и профессору Александру Гасникову.


\begin{quote}\small
\textbf{Статус: исследовательская ветка, всё приложение находится на стадии проверки.}
Эта ветка продолжает изучать следствия найденных выпрямлений и их связи
с проблемой вложений Абъянкара--Моха, группами Кремоны, бирациональной
и аффинной геометрией. Несмотря на предварительный статус, авторы считают
полезным поделиться аргументами и вычислениями. Авторы надеются, что они
вырастут в полноценную исследовательскую программу, которая поможет
продвинуться в понимании этих областей, и приглашают всех заинтересованных
к проверке, обсуждению, исправлению и дальнейшему исследованию.

Материал подготовлен в ходе работы с ИИ. Внутренние повторные проверки и
альтернативные реализации не являются внешним независимым рецензированием.
Названия <<теорема>> и <<доказательство>> обозначают предлагаемые математические
аргументы; они не отменяют общего статуса проверки. Публикационный приоритет
не заявляется. Приложение не заменяет исходный препринт и Appendix 1.
\end{quote}

\paragraph{Источник и границы обновления.}
Зафиксирован репозиторий
\href{https://github.com/qft98romanov-hash/abhyankar-moh-trefoil/tree/b954f3609e900761fddc9c8ff03d8301e9604cab}{\nolinkurl{qft98romanov-hash/abhyankar-moh-trefoil}},
commit \nolinkurl{b954f3609e900761fddc9c8ff03d8301e9604cab}.
Исходные формулы и ранее опубликованные приложения сохранены.
Здесь слово <<закрыт>> означает логическое следствие приведённых формул
и цитируемых теорем внутри данной версии, а не заявление об общем решении
проблемы или о принятии результата математическим сообществом.

\tableofcontents
\section{Что установлено, какие кандидаты исключены и что остаётся открытым}
\label{sec:status-literature}

\paragraph{Статус исследовательской ветки.}
Все утверждения этого приложения представлены как рабочие доказательства и
вычислительные сертификаты, находящиеся на стадии проверки. Внутренняя
проверка тождеств и повторное применение опубликованных критериев не равны
внешнему независимому рецензированию. Ниже слово \emph{следствие} обозначает
логический вывод из выписанных формул и указанных теорем; оно не означает
заявления о литературном приоритете.

\begin{center}
\small
\begin{tabular}{@{}p{0.30\linewidth}p{0.64\linewidth}@{}}
\toprule
Вопрос & Результат и его точная граница\\
\midrule
Конкретная кривая Шастри &
Полиномиальный выпрямитель над $K=\Q(\sqrt{-3})$ исключает эту
параметризацию из кандидатов на комплексную невыпрямляемость. Он не
исключает существования других невыпрямляемых пространственных прямых.\\[4pt]
Гармоническая модель и $C_N$ &
Выпрямляются именно указанные параметризации. Гармоническая модель
использует $D_5=t^5-5t^3+5t$; замена $5t$ на $4t$ не покрывается
этим утверждением.\\[4pt]
Ручность построенных карт &
Дикость $\Phi_T,\Phi_H,\Phi_N$ следует из критериев относительных
координат. Отсутствие \emph{любого} ручного выпрямителя тех же кривых
не доказано.\\[4pt]
Стабилизация &
$\Phi_\ast\times\id_{\A^3}\in\Tame_6(K)$ следует из известной
теоремы о стабильной ручности. Минимальное число дополнительных
переменных не найдено.\\[4pt]
Кремонова подгруппа &
Построенные карты принадлежат $\mathcal C_3(K)$.
Это не описание всей $\Cr_3(K)$ и не первый пример дикого элемента
в $\mathcal C_3(K)$.\\[4pt]
Поле определения &
Для Шастри минимальная степень подходящего числового поля над $\Q$
равна $2$. Единственность подходящего квадратичного поля и эффективная
классификация полей выпрямления не установлены.\\
\bottomrule
\end{tabular}
\end{center}

\paragraph{Исторический кандидат и общая проблема.}
В работе Шастри \cite[\S1 и \S3, Example~1]{Shastri1992} вещественные
узловые параметризации предложены как кандидаты для исследования
комплексной невыпрямляемости. Поэтому наличие полного полиномиального
автоморфизма для $\gamma_T$ решает вопрос о \emph{данной кривой}
положительно. Историческая гипотеза о существовании хотя бы одного
невыпрямляемого вложения $\A^1_{\C}\hookrightarrow\A^3_{\C}$ этим не
опровергается. Равным образом не доказана противоположная универсальная
гипотеза о выпрямляемости всех таких вложений. Мы не приписываем статье
Шастри отдельную нумерованную гипотезу, которой там не сформулировано.

\paragraph{Связанные гипотезы из первого приложения.}
Нужно отдельно учитывать уже записанные в первом приложении следствия
для кандидатов ван ден Эссена--ван Россума
\cite[\S5, Conjectures~5.1--5.3]{VdER2002}.\footnote{Нумерация и
сопоставление объектов здесь унаследованы из Appendix 1. При подготовке
этой версии издательская запись была проверена, но повторная загрузка
полного текста AMS оказалась недоступна. Алгебраические выводы ниже
проверены непосредственно; это ограничение библиографического аудита
зафиксировано также в \texttt{audit/claim\_status.md}.}
Это продолжение исходного выпрямления, а не новый результат о группе
Кремоны. При $\gamma_T=(f,g,h)$ рассматривается
\[
 D=t\partial_u+f'(u)\partial_x+g'(u)\partial_y+h'(u)\partial_z
 \quad\text{на }\C[t,u,x,y,z].
\]
Гипотеза~5.1 (в последующей литературе также \emph{Conjecture E})
предсказывала неполиномиальность $\ker D$. Из выпрямителя
$\Phi_T=(P,Q,R)$ с обратным $\Psi$ получаются, при
$V=\gamma_T(u)-t(x,y,z)$,
\[
 p=P(V),\qquad q=Q(V)/t,\qquad r=R(V)/t,\qquad
 s=(u-P(V))/t.
\]
Все деления точны, поскольку $\Phi_T\gamma_T=(u,0,0)$. Обратная замена
равна
\[
 u=p+ts,\qquad
 (x,y,z)=\frac{\gamma_T(p+ts)-\Psi(p,tq,tr)}{t};
\]
она также полиномиальна. Следовательно,
\[
 \C[t,u,x,y,z]=\C[t,p,q,r,s],\qquad
 D=\partial_s,\qquad \ker D=\C[t,p,q,r].
\]
Это даёт отрицание именно указанного кандидата~5.1. Если $\sigma$
обозначает исходный слайс из того же построения, то
$\sigma-s\in\ker D$; отсюда
$\C[t,u,x,y,z]/(\sigma)\simeq\C^{[4]}$.
Также при $I_{\gamma_T}=(Q,R)$ имеем
\[
 \C[x,y,z][t,I_{\gamma_T}/t]
 =\C[t,P,Q/t,R/t]\simeq\C^{[4]}.
\]
Это соответствующие полиномиальные описания объектов~5.2 и~5.3;
подробные формулы, сопоставление обозначений и сертификаты находятся
в первом приложении репозитория. Общие проблемы сокращения,
вложений и линеаризации действий из этих частных вычислений не решаются.
Отозванный препринт Оды \cite{OdaWithdrawn} не используется как
доказательство: его запись arXiv содержит отзыв от 9 мая 2020 года
с указанием ошибки.

\paragraph{Предшественники программы.}
Связь задачи вложений с кремоновой эквивалентностью явно обсуждается
у Меллы--Поластри \cite[введение и Theorem~1]{MP2009}.
Относительные координаты, факторизации над полем частных и перенос
изучаются у Эдо--Венеро \cite{EV2001}.
Наиболее близкий предшественник конструктивной части --- Льюис
\cite[Lemma~1, Theorems~2 и~12]{Lewis2024}: координата, ограничивающаяся
в параметр, строится с использованием остаточных координат и
исправления рациональных преобразований. В этой же работе сопоставлены
предшествующие результаты Крайгеро, Бхатвадекара--Роя и Куроды.
Поэтому общий приём ``локализация, затем устранение знаменателей''
не заявляется здесь как новая стратегия.

Наличие диких автоморфизмов в трёх переменных давно установлено;
перенос дикости относительной координаты в три переменные даёт
\cite{UY2004}. Блан--Хеден уже доказали
$\operatorname{Nagata}\in\mathcal C_3$
\cite[Proposition~6.8]{BH2015}. Стабильная ручность используемого
класса следует из \cite[Corollary~4.9]{BEW2012}.
Самостоятельное содержание данного приложения состоит в применении
этих результатов к \emph{конкретным} выпрямителям, проверке их
гипотез и явном описании дальнейшей программы.

\paragraph{Нерешённые вопросы и границы поиска.}
Открытыми в нашей работе остаются полнота относительного метода,
существование ручного выпрямителя для тех же кривых, минимальные степени
выпрямителей, спуск в другие вполне мнимые числовые поля и общий случай
$\A^1\hookrightarrow\A^3$ над $\C$.
Конечные вычисления для $N\leq N_{\max}$ служат проверкой реализации;
утверждения для всех $N$ опираются на индуктивные доказательства.
Неуспех ограниченного поиска не доказывает невыпрямляемость.
Целевой литературный поиск по состоянию на 28 октября 2026 года не
выявил более ранней публикации с нашим полным выпрямителем Шастри,
однако отсутствие найденной работы не устанавливает приоритет.

\section{Постановка и программа исследования}\label{sec:program}

Во всём приложении \emph{выпрямлением} называется автоморфизм всего
аффинного пространства, переводящий заданную параметризованную кривую
в стандартную прямую. Выбор поля и требование полиномиальности обратного
отображения входят в постановку. Бирациональное выпрямление является
более слабым условием.

\subsection{Координатный подъём параметра}

Пусть $k$ --- поле, $R=k[x,y,z]$ и
\[
 \gamma(t)=(p(t),q(t),r(t)),\qquad p,q,r\in k[t].
\]
Параметризация задаёт замкнутое вложение $\A^1_k\hookrightarrow\A^3_k$
тогда и только тогда, когда
\begin{equation}\label{eq:closed-embedding}
 k[p,q,r]=k[t].
\end{equation}
Это эквивалентно сюръективности гомоморфизма ограничения
$\gamma^*:R\to k[t]$. В частности, существует полином $H_0\in R$
с $H_0\gamma=t$. Однако такой полином ещё не обязан быть
\emph{координатой}, то есть компонентой полиномиального автоморфизма.

\begin{lemma}[Координатный подъём]\label{lem:coordinate-lift}
Вложение $\gamma$ допускает выпрямитель $\Phi\in\Aut_k(\A^3)$,
$\Phi\gamma(t)=(t,0,0)$, тогда и только тогда, когда существует
координата $H\in R$, удовлетворяющая $H\gamma=t$.
\end{lemma}
\begin{proof}
Необходимость даёт первая компонента $\Phi$. Обратно, дополним $H$
до автоморфизма $(H,U,V)$. На кривой его значение равно
$(t,u(t),v(t))$. Применим затем автоморфизм
\[
 (T,Y,Z)\longmapsto(T,Y-u(T),Z-v(T)).
\]
Полученная композиция выпрямляет $\gamma$.
\end{proof}
Критерий известен; он явно используется в \cite[лемма~1]{Lewis2024}.
Его конструктивное содержание состоит в выборе координаты среди всех
подъёмов $H_0+\ker\gamma^*$.

\subsection{Что даёт двумерный случай}

В характеристике нуль теорема Абъянкара--Моха--Судзуки утверждает,
что всякое замкнутое вложение $\A^1_k\hookrightarrow\A^2_k$
выпрямляется \cite{AM1975,Suzuki1974}. Для непараметризованных образов
это означает транзитивность действия $G_2(k)=\Aut_k(\A^2)$.
Положив $L_2=\{y=0\}$, получаем описание множества таких образов:
\[
 \{\text{образы замкнутых вложений }\A^1_k\hookrightarrow\A^2_k\}
 \simeq G_2(k)/\operatorname{Stab}(L_2).
\]
Здесь используются левые смежные классы $g\operatorname{Stab}(L_2)$, причём
\begin{equation}\label{eq:plane-stabilizer}
 \operatorname{Stab}(L_2)=
 \{(x,y)\mapsto(ax+P(y),by):a,b\in k^*,\ P\in k[y]\}.
\end{equation}
Действительно, сохранение идеала $(y)$ даёт образ $y$, равный $by$.
После масштабирования остаётся автоморфизм $k[y][x]$, поэтому образ $x$
имеет вид $ax+P(y)$ с $a\in k^*$. Обратное включение непосредственно.

Следует различать три классических утверждения. Теорема
Абъянкара--Моха--Судзуки описывает орбиту вложенных прямых. Теорема
Юнга--ван дер Кулка описывает $G_2(k)$ через аффинные и треугольные
автоморфизмы и справедлива над любым полем
\cite[теорема~1]{EV2001}. Теорема Нётера--Кастельнуово утверждает
порождение $\Cr_2(\C)$ проективными автоморфизмами и стандартной
квадратичной инволюцией; см. обсуждение в \cite[\S1]{BH2015}.
Описание орбиты прямой само по себе не является описанием всей группы
Кремоны. В положительной характеристике общая плоская теорема о
выпрямлении неверна, поэтому перенос её заключения на произвольное
поле без оговорок недопустим.

\subsection{Бирациональное выпрямление и точная групповая постановка}

Выбор стандартной аффинной карты даёт включение
\[
 G=\Aut_k(\A^3)\subset B=\Cr_3(k)=\operatorname{Bir}_k(\mathbb P^3).
\]
Для $b\in B$ имеется точный критерий:
\begin{equation}\label{eq:ring-preservation}
 b\in G\quad\Longleftrightarrow\quad b^*R=R.
\end{equation}
Действительно, равенство означает полиномиальность всех компонент $b$
и $b^{-1}$. В терминах порядков вдоль неприводимых полиномов $P\in R$
это условие записывается как
\[
 v_P(b_i)\geq0,\qquad v_P((b^{-1})_i)\geq0
 \quad\text{для всех }P\text{ и }i.
\]
Используется факториальность $R$. Полюса на гиперплоскости бесконечности
при этом допустимы. Проверка одного только якобиана не заменяет
условия~\eqref{eq:ring-preservation}.

\begin{proposition}\label{prop:birational-any-field}
Над любым полем $k$ всякое вложение~\eqref{eq:closed-embedding}
имеет бирациональный выпрямитель над $k$. Его и обратное можно выбрать
регулярными около общих точек соответствующих кривых. Рациональный
якобиан выпрямителя можно сделать равным $1$.
\end{proposition}
\begin{proof}
Дифференцирование равенства $H_0(p,q,r)=t$ даёт
\[
 (H_0)_x(p,q,r)p'+(H_0)_y(p,q,r)q'+(H_0)_z(p,q,r)r'=1.
\]
Поэтому после циклической перестановки координат можно считать
$p'\ne0$. Расширение $k(t)/k(p)$ конечно и сепарабельно, и
$k(t)=k(p)(q,r)$. По теореме о примитивном элементе можно выбрать
$\lambda(p)\in k[p]$ так, чтобы
$k(p,q+\lambda(p)r)=k(t)$. В самом деле, плохих коэффициентов
$\lambda\in k(p)$ лишь конечное число, а $k[p]$ бесконечно даже
при конечном $k$. Это следует также непосредственно из сравнения
различных вложений сепарабельного расширения: каждое неравное
сопряжение исключает не более одного значения $\lambda$.

Применим сдвиг $(x,y,z)\mapsto(x,y+\lambda(x)z,z)$ и снова обозначим
вторую компоненту кривой через $q$. Существует
$\tau\in k(x,y)$ с $\tau(p,q)=t$; знаменатель выбранного выражения
не равен нулю тождественно на кривой. Положим
\begin{equation}\label{eq:birational-construction}
 s=z-r(\tau(x,y))+\tau(x,y),\qquad
 b=(s,x-p(s),y-q(s)).
\end{equation}
Тогда $b\gamma=(t,0,0)$. Для входа $(T,U,V)$ обратное задаётся формулами
\begin{equation}\label{eq:birational-inverse}
 \begin{split}
 X&=U+p(T),\qquad Y=V+q(T),\\
 b^{-1}(T,U,V)&=(X,Y,T+r(\tau(X,Y))-\tau(X,Y)).
 \end{split}
\end{equation}
Обе композиции тождественны непосредственной подстановкой. На прямой
$(T,0,0)$ знаменатель становится тем же ненулевым элементом $k(T)$,
поэтому обратное также определено в общей точке. Все использованные
сдвиги имеют якобиан $1$, как и циклические перестановки.
\end{proof}

Над алгебраически замкнутым полем бирациональная эквивалентность
проективного замыкания нашей кривой и прямой является частным случаем
теоремы Меллы--Поластри: два бирациональных вложения $d$-мерного
многообразия в $\mathbb P^n$ кремоново эквивалентны при $n\geq d+2$
\cite[теорема~1]{MP2009}. Приведённое элементарное доказательство
указывает карту над исходным полем. Ни одно из этих утверждений
не гарантирует регулярность во всех точках аффинного пространства.

Для Шастри можно явно взять
\[
 \tau(x,y)=\frac{x(y+1)}{x^2-y-3},\qquad
 s=z-\tau^5+11\tau,\qquad
 b=(s,x-s^3+3s,y-s^4+4s^2).
\]
Равенство $p(q+1)=t(p^2-q-3)$ при $p=t^3-3t$, $q=t^4-4t^2$
проверяет восстановление параметра. Это бирациональный выпрямитель
над $\Q$, допускающий полюса; обе его композиции с явным обратным
проверяются в новом вычислительном пакете. Его существование совместимо
с невозможностью полиномиального выпрямления над $\Q$, обсуждаемой ниже.

Пусть $L=\{y=z=0\}$. Обозначим через $S_L$ подгруппу $B$,
состоящую из бирациональных карт, задающих изоморфизмы между некоторыми
окрестностями общей точки $L$ и действующих на $L$ тождественно.
Таким образом, условие определённости предъявляется и к обратной карте.

\begin{proposition}[Задача пересечения смежного класса]\label{prop:coset}
Если $b$ --- бирациональный выпрямитель из
предложения~\ref{prop:birational-any-field}, то
\begin{equation}\label{eq:coset}
 \gamma\text{ полиномиально выпрямляется}
 \quad\Longleftrightarrow\quad S_Lb\cap G\ne\varnothing.
\end{equation}
\end{proposition}
\begin{proof}
Для полиномиального выпрямителя $\Phi$ карта $h=\Phi b^{-1}$
принадлежит $S_L$, и $\Phi=hb$. Обратно, если $hb\in G$ при
$h\in S_L$, то $(hb)\gamma(t)=(t,0,0)$ как рациональное равенство.
Поскольку обе стороны полиномиальны, равенство справедливо всюду.
\end{proof}
Таким образом, исходная задача переформулируется как поиск поправки
$h$, устраняющей аффинные полюса $b$ и его обратного при сохранении
параметризованной прямой. Формула~\eqref{eq:coset} сама по себе не
является алгоритмом и не доказывает существование нужного пересечения.

\subsection{Относительная подгруппа как рабочий механизм}

Стабилизатор рациональной функции $x$ и его полиномиальная часть равны
\[
 J_x=\{b\in B:b^*x=x\}\simeq\Cr_2(k(x)),\qquad
 G_x=G\cap J_x=\Aut_{k[x]}k[x,y,z].
\]
Элемент $G_x$ становится ручным автоморфизмом двух переменных над
$k(x)$, но его факторизация над этим полем может иметь знаменатели.
Их устранение над $k[x]$ --- содержательная часть конструкции.
Этот подход имеет прямых предшественников в теории переноса координат
и остаточных координат \cite{EV2001,Lewis2024}.

Если $A=(x,H,W)\in G_x$, $H\gamma=t$ и $W\gamma=w(t)$, то
\begin{equation}\label{eq:relative-rectifier}
 \Phi=(H,x-p(H),W-w(H))
\end{equation}
выпрямляет $\gamma$. В самой группе $G_x$ естественная целевая
параметризация имеет вид $(p(t),t,0)$: фиксирующие $x$ карты не меняют
первую компоненту $p(t)$. Последующий внешний ручной автоморфизм
переводит её в $(t,0,0)$.

\begin{remark}\label{rem:relative-not-necessary}
Для заранее фиксированного $x$ существование относительной координаты
$H$ с $H\gamma=t$ является достаточным условием выпрямляемости.
Его необходимость для произвольного выпрямляемого вложения здесь не
доказана. Разрешение произвольной предварительной полиномиальной
замены без дополнительных ограничений сделало бы утверждение о
полноте круговым: такой заменой мог бы оказаться искомый выпрямитель.
\end{remark}

\subsection{Поле определения и препятствие к спуску}

Для кривой Шастри
\[
 \gamma_T=(t^3-3t,t^4-4t^2,t^5-10t)
\]
вещественный образ задаёт трилистник \cite[\S3, пример~1]{Shastri1992}.
Полиномиальный автоморфизм $\mathbb R^3$ является собственным
диффеоморфизмом и продолжается до гомеоморфизма одноточечной
компактификации $S^3$. Он не переводит нетривиальный узел в тривиальный.
Поэтому выпрямление над $\mathbb R$, а также над $\Q$, невозможно.
Тот же аргумент применим к любому числовому полю с вещественным
вложением: достаточно применить его к коэффициентам предполагаемой карты.

В сочетании с явным выпрямителем над $K=\Q(\sqrt{-3})$ это даёт
\[
 \min\{[F:\Q]: F\text{ --- числовое поле, над которым }
 \gamma_T\text{ выпрямляется}\}=2.
\]
Подходящее числовое поле должно быть вполне мнимым. Единственность
$K$ среди квадратичных полей не установлена; отсутствие в поле корня
$a^2-a+1$ исключает непосредственную специализацию нашей формулы,
но не другую конструкцию.

Полиномиальный спуск допускает точную групповую запись. Пусть
$\alpha\in\Aut_K(\A^3)$ выпрямляет $\gamma_T$,
$\Gamma=\operatorname{Gal}(K/\Q)$ и $H_L(K)$ --- поточечный
\emph{полиномиальный} стабилизатор $L$. Тогда
\[
 c_\sigma=\alpha\,\sigma(\alpha)^{-1}\in H_L(K),\qquad
 c_{\sigma\tau}=c_\sigma\,\sigma(c_\tau).
\]
Выпрямитель над $\Q$ существует тогда и только тогда, когда найдётся
$h\in H_L(K)$ с $c_\sigma=h^{-1}\sigma(h)$ для всех $\sigma$.
Действительно, это равносильно определённости $h\alpha$ и его
полиномиального обратного над $\Q$. Для Шастри такой $h$ не существует.
После замены полиномиального стабилизатора на $S_L$ препятствие исчезает:
при бирациональном выпрямителе $b$ над $\Q$ подходит $h=b\alpha^{-1}$.
Следовательно, переход к полной группе Кремоны теряет в том числе
информацию о полиномиальном спуске. Здесь используется лишь
элементарное равенство кокциклов; классификация некоммутативного
множества $H^1(\Gamma,H_L(K))$ не заявляется.

\subsection{Что остаётся программой, а не результатом}

\begin{enumerate}
\item \textbf{Полнота переносов.} Описать класс параметризаций,
для которого выбор центров модификаций и исключительных дивизоров
всегда завершается относительной координатой. Конечные вычисления
для отдельных индексов не доказывают такую полноту.
\item \textbf{Контролируемые изменения проекции.} Выяснить, какие
заранее заданные аффинные или ограниченные полиномиальные замены
позволяют расширить применимость $G_x$ без круговой постановки.
\item \textbf{Ручная орбита самой кривой.} Дикость построенного
выпрямителя не исключает другого ручного выпрямителя. Требуется
исследовать все координаты в $H+\ker\gamma^*$, а не только выбранную $H$.
\item \textbf{Арифметика.} Определить другие поля выпрямления,
например проверить существование выпрямителя Шастри над $\Q(i)$.
Отсутствие корня в конкретной формуле и препятствие к существованию
любого автоморфизма --- разные утверждения.
\end{enumerate}
Эти цели не требуют предварительного полного описания $\Cr_3(\C)$.
Они исследуют её определённые подгруппы, смежные классы и условия
регулярности на фиксированной аффинной модели.

\section{Конструкции, используемые в групповых выводах}\label{sec:constructions}

\subsection{Поле и базовая координатная пара}
Положим
\[
 K=\Q[a]/(a^2-a+1),\quad b=1-a,\quad ab=1,\quad b^2=-a.
\]
Чебышёвские полиномы в целочисленной нормировке определяются рекурсией
$D_0=2$, $D_1=t$, $D_{n+1}=tD_n-D_{n-1}$. В характеристике нуль
$D_n(t)=2T_n(t/2)$. В частности,
\[
 D_3=t^3-3t,\quad D_4=t^4-4t^2+2,\quad D_5=t^5-5t^3+5t,
\quad D_7=t^7-7t^5+14t^3-7t.
\]
Рассматриваем конкретные параметризации
\[
 \gamma_T=(D_3,D_4-2,t^5-10t),\qquad
 \gamma_H=(D_3,D_5,D_7).
\]
Во второй компоненте $\gamma_H$ коэффициент при $t$ равен $5$.
Параметризация с коэффициентом $4$ данным результатом не охватывается.

Аффинная карта $M=(x,y+2,z-5x)$ переводит $\gamma_T$ в
$\gamma_0=(D_3,D_4,D_5)$. В нормализованных координатах введём
\begin{align*}
 u&=axy+z,&v&=-y,&m&=bx^5,&s&=bx,\\
 \ell&=-1+ax^2-ax^4,&
 g&=-2x^5+(6-2a)x^3-2bx,\\
 P_x(U)&=sU^2+\ell U+g,&H&=P_x(u)+mv.
\end{align*}
В кольце $K[x,T]/(x^5)$ полином $P_x$ обратим. Используем представитель
\begin{align}
 F_x(T)={}&-T+x(bT^2+2a-2)+x^2(2aT^3-5aT)\notag\\
 &+x^3(-5T^4+15T^2-2a)\notag\\
 &+x^4\bigl(14bT^5+50(a-1)T^3+(29-24a)T\bigr).
 \label{eq:inverse-mod-five}
\end{align}
Прямое вычисление даёт обе конгруэнции
\[
 P_x(F_x(T))\equiv T\pmod{x^5},\qquad
 F_x(P_x(U))\equiv U\pmod{x^5}.
\]
Поэтому
\[
 W=\frac{u-F_x(H)}m,\qquad
 \chi_x(T)=\frac{T-P_x(F_x(T))}m
\]
являются полиномами. Для $A_0=(x,H,W)$ обратная карта имеет вид
\begin{align}
 U&=F_x(T)+mV,\notag\\
 V_0&=\chi_x(T)-(\ell+2sF_x(T))V-smV^2,\notag\\
 A_0^{-1}(x,T,V)&=(x,-V_0,U+axV_0).\label{eq:base-inverse}
\end{align}
Действительно, раскрытие квадрата даёт $P_x(U)+mV_0=T$ и
$(U-F_x(T))/m=V$. Для обратной композиции сначала восстанавливаются
$u,v$, затем $y,z$. Это полиномиальное доказательство обратимости;
постоянство якобиана само по себе здесь не используется как критерий.

Подстановки в кривую дают
\[
 H(D_3,D_4,D_5)=t,\qquad W(D_3,D_4,D_5)=8-14t^2.
\]
Следовательно, при $h_T=H\circ M$, $w_T=W\circ M$ получаем
\begin{equation}\label{eq:trefoil-rectifier}
 \Phi_T=(h_T,\ x-h_T^3+3h_T,\ w_T-8+14h_T^2),
 \qquad \Phi_T\gamma_T=(t,0,0).
\end{equation}
Обратное строится из \eqref{eq:base-inverse}, $M^{-1}$ и обратных
внешних сдвигов. Степени компонент данной карты равны $(5,15,24)$,
а якобиан равен $1$. Оптимальность степеней не заявляется.
Исходная реализация находится в \texttt{verify.py} репозитория;
новый проверяющий код отдельно воспроизводит компактные сертификаты.

\subsection{Перенос с контролем обеих обратных карт}
Следующий механизм родственен остаточным координатам
\cite{BD1993,EV2001,Lewis2024}; здесь записаны формулы, необходимые
для воспроизводимости.

\begin{proposition}[Достаточный критерий переноса]\label{prop:transfer-app2}
Пусть $A=(x,H,W)$ и $A^{-1}=(x,F,G)$ -- полиномиальные взаимно
обратные карты над $k$. Пусть $0\ne m\in k[x]$, $p\in k[x,y]$ и
\[
 \beta=(x,p(x,y)+m(x)z,y).
\]
Если $c(y)=\overline{H(x,p(x,y),y)}$ -- автоморфизм
$(k[x]/(m))[y]$, то $H\circ\beta$ дополняется до относительного
полиномиального автоморфизма.
\end{proposition}
\begin{proof}
Пусть $\psi$ -- обратный к $c$ в кольце $B[T]$, где $B=k[x]/(m)$.
Выберем подъём полинома
\[
 \varphi(x,T)\equiv W(x,p(x,\psi(T)),\psi(T))\pmod m.
\]
Тогда
\begin{equation}\label{eq:transfer-forward}
 A_{\rm new}=
 \left(x,H\circ\beta,
 \frac{W\circ\beta-\varphi(x,H\circ\beta)}m\right).
\end{equation}
Числитель последней компоненты обращается в нуль по модулю $m$.
Для входа $(x,T,V)$ положим
$S=mV+\varphi(x,T)$, $Y=F(x,T,S)$, $Z=G(x,T,S)$. Обратная карта:
\begin{equation}\label{eq:transfer-inverse}
 A_{\rm new}^{-1}(x,T,V)=\left(x,Z,\frac{Y-p(x,Z)}m\right).
\end{equation}
В $B[T]$ из обратимости $A$ получаем
$Y=p(x,\psi(T))$, $Z=\psi(T)$; деление снова точное.
После обращения $m$ формулы представляют композиции явно обратимых
карт. Инъективность локализации переносит обе тождественные композиции
в $k[x,y,z]$, включая все слои $m=0$.
\end{proof}

Условие на $c=c_0+c_1y+\sum_{j\ge2}c_jy^j$ эквивалентно обратимости
$c_1$ и нильпотентности всех $c_j$, $j\ge2$, в $B$.
Если $J=(c_j:j\ge2)$, $J^e=0$, обратный находится конечной итерацией
\[
 \psi_0=c_1^{-1}(T-c_0),\qquad
 \psi_{i+1}=c_1^{-1}\left(T-c_0-\sum_{j\ge2}c_j\psi_i^j\right).
\]
Разности последовательных итераций получают всё более высокие степени
$J$; после $e$ шагов достигается неподвижная точка. Таким образом,
вычисления должны вестись по полному $m$, а не только по его радикалу.

\subsection{Гармоническая модель и семейство $C_N$}\label{subsec:family-app2}
Тождество
\[
 D_3D_7-D_5^2+2=D_4
\]
показывает, что $\beta_1=(x,xz-y^2+2,y)$ переводит $\gamma_H$ в
$\gamma_0$. Поскольку $H(0,Y,Z)=-Z$, центр переноса равен $-y$
по модулю $x$, и предложение~\ref{prop:transfer-app2} применимо.
Можно взять поправку
\[
 \varphi_1(T)=42(a-1)T^6+175(1-a)T^4+(147a-161)T^2+8.
\]
Полагая $H_1=H\circ\beta_1$,
$W_1=(W\circ\beta_1-\varphi_1(H_1))/x$, имеем
\[
 W_1\gamma_H=b(42t^3-49t),\quad
 \Phi_H=(H_1,x-H_1^3+3H_1,W_1-b(42H_1^3-49H_1)).
\]
Степени компонент равны $(7,21,41)$. Обратное задаётся
\eqref{eq:transfer-inverse} и внешними обратными сдвигами.

Теперь пусть $N\ge1$. В $\Z[x,y]/(x^N)$ определим
\[
 \nu_0=-y,\qquad \nu_{i+1}=-y+x(\nu_i^2-2),\quad 0\le i<N,
 \qquad p_N\equiv2+x\nu_N-\nu_N^2\pmod{x^N},
\]
выбирая остаток $p_N$ степени по $x$ меньше $N$. Имеем
$p_N\equiv2-y^2\pmod x$. Положим
\begin{equation}\label{eq:family-CN}
 C_N(t)=\left(D_3,D_5,
 \frac{D_4-p_N(D_3,D_5)}{D_3^N}\right),\qquad
 \beta_N=(x,p_N+x^Nz,y).
\end{equation}
Деление в параметризации точное. Действительно,
$D_5=D_3(t^2-2)-t$ и $D_4=D_3t-t^2+2$.
Ошибка итерации $\nu_i(D_3,D_5)-t$ при каждом шаге получает дополнительный
множитель $D_3$. Следовательно, $\nu_N(D_3,D_5)\equiv t\pmod{D_3^N}$,
откуда и следует делимость.

Центр $H(x,p_N,y)$ равен $-y$ по модулю $x$. Поэтому по модулю $x^N$
его линейный коэффициент обратим, а все старшие нильпотентны.
Предложение~\ref{prop:transfer-app2} строит $A_N=(x,H_N,W_N)$ с
$H_N=H\circ\beta_N$ и $H_NC_N=t$. При $w_N=W_NC_N$ получаем
\begin{equation}\label{eq:family-rectifier}
 \Phi_N=(H_N,x-D_3(H_N),W_N-w_N(H_N)),\qquad
 \Phi_NC_N=(t,0,0).
\end{equation}
Тем самым $C_N$ действительно задаёт замкнутое вложение.
При $N=1$ можно выбрать указанный выше перенос и получить
$C_1=\gamma_H$, $\Phi_1=\Phi_H$. Различные представители поправки
могут менять компаньон ручным сдвигом; все последующие утверждения
относятся к построенным таким образом картам.

\subsection{Интегральность формул}
Все базовые коэффициенты лежат в
$\mathcal O=\Z[a]/(a^2-a+1)$. Конгруэнции для $P_x,F_x$ верны
уже над $\mathcal O$, а $b^{-1}=a$. Деление на $bx^5$ состоит из
умножения на $a$ и точного деления на моном $x^5$.
Рекурсии для $p_N$ и обратного центра не вводят числовых знаменателей:
линейный коэффициент центра равен $-1$ по модулю $x$ и обращается
конечным геометрическим рядом по модулю $x^N$.
Поэтому полиномиальные карты и их обратные специализируются в любое поле,
содержащее корень $a^2-a+1$, включая характеристики $2$ и $3$.
Это утверждение о формулах; доказательства дикости ниже относятся к
характеристике нуль. Степени после специализации могут уменьшаться.
Для $\mathbb F_q$ корень существует тогда и только тогда, когда
$\operatorname{char}\mathbb F_q=3$ или $q\equiv1\pmod3$:
в характеристике, отличной от $3$, элемент $-a$ является нетривиальным
кубическим корнем единицы. Отсутствие корня не исключает иных выпрямителей.

\section{Дикость построенных выпрямителей}\label{sec:group-results}

В этом разделе $K=\Q[a]/(a^2-a+1)$, $b=1-a$, $R=K[x]$.
Используются равенства $ab=1$ и $b^2=-a$.
Ручная группа $\Tame_3(K)$ порождена аффинными и элементарными
полиномиальными автоморфизмами. Координата называется ручной, если она
входит в ручной автоморфизм, и дикой в противном случае.
Относительная ручность над $R$ означает ручность в двух переменных над
кольцом $R$; это условие первоначально сильнее ручности над полем
коэффициентов $K$. Приведённые ниже выводы из опубликованных теорем
являются рабочими доказательствами этой исследовательской ветки и
предлагаются для дальнейшей проверки.

\begin{proposition}[Выпрямитель Шастри]\label{prop:trefoil-wild}
Координата $h_T$ и построенный выпрямитель $\Phi_T$ являются дикими
над $K$.
\end{proposition}
\begin{proof}
В относительных координатах $(u,v)=(axy+z,-y)$ базовая координата равна
$H=mv+P_x(u)$, где $m=bx^5$, а коэффициент при $u^2$ равен $bx$.
Критерий Эдо--Венеро для координат длины один
\cite[предложение~2]{EV2001} требует для ручности делимости всех
нелинейных коэффициентов на $m$. Но $bx^5\nmid bx$ в $R$.
Замена $(y,z)\mapsto(u,v)$ ручная над $R$, поэтому $H$ дикая над $R$.

По теореме Умирбаева--Ю дикая $K[x]$-координата остаётся дикой
координатой $K[x,y,z]$ в характеристике нуль
\cite[главная теорема]{UY2004}.
Предварительная замена $M=(x,y+2,z-5x)$ ручная, поэтому $h_T=H\circ M$
тоже дикая. Ручный автоморфизм не может иметь дикую компоненту;
следовательно, $\Phi_T$ дикий.
\end{proof}

\subsection{Внутренняя координата рациональной длины два}

Для следующей леммы $Y,Z$ обозначают две переменные над $R$.
Запись
\[
 F=p_2Y+Q_2(p_1Z+Q_1(Y))
\]
называем приведённой в смысле Эдо, если $p_1$ монический,
$Q_1\in R[Y]$, $Q_1(0)=Q_2(0)=0$,
$\gcd(p_1,\text{коэффициенты }Q_1)=1$;
здесь $p_1,p_2\ne0$, $p_2\in K(x)$, $Q_2\in K(x)[Y]$.
Мы используем случай $\deg Q_1\ge2$, $\deg Q_2\ge2$.

\begin{lemma}[Следствие классификации Эдо]\label{lem:inner-wild}
Пусть $F\in R[Y,Z]$ является координатой и имеет указанную
приведённую запись. Если $F$ ручная над $R$, то
$p_1Z+Q_1(Y)$ также ручная над $R$.
\end{lemma}
\begin{proof}
Здесь композиции $\sigma\tau$ понимаются как композиции
$R$-алгебраических автоморфизмов: $(\sigma\tau)(f)=\sigma(\tau(f))$.
Теорема~8(1) Эдо даёт $F=\sigma\tau(Y)$, где обе компоненты каждого
из $\sigma,\tau$ имеют длину один; такие автоморфизмы ручные.
Как в доказательстве теоремы~8(3), нормируем
$\sigma(Y)=p_3Z+Q_3(Y)$, где $p_3$ монический и $Q_3(0)=0$.
Нормировка аффинная и сохраняет ручность.
Формула компаньона Рассела--Сатайе имеет вид
\[
 \sigma(Z)=\frac{\varepsilon}{p_3}
 \bigl(V(p_3Z+Q_3(Y))-Y\bigr),\qquad
 \varepsilon\in R^\times,\quad V\in R[Y].
\]
Если $\tau(Y)=p_4Z+Q_4(Y)$, то
\[
 F=-\frac{\varepsilon p_4}{p_3}Y+
 \left(\frac{\varepsilon p_4}{p_3}V+Q_4\right)
 (p_3Z+Q_3(Y)).
\]
После удаления постоянного слагаемого это приведённая запись того же
полинома. Её единственность, установленная в предложении~2 Эдо,
даёт $p_3=p_1$, $Q_3=Q_1$. Поэтому внутренний полином является
компонентой ручной $\sigma$; см.\ также ту же нормировку в
доказательстве теоремы~8(3) \cite{Edo2013}.
\end{proof}

\begin{proposition}[Гармоническая модель и все построенные $C_N$]
\label{prop:family-wild}
Для каждого $N\ge1$ координата $H_N$ и построенный выпрямитель
$\Phi_N$ дикие над $K$. В частности, это верно для гармонической
модели $C_1=\gamma_H$.
\end{proposition}
\begin{proof}
Положим
\[
 Q_x(U)=bxU^2+(-1+ax^2)U+g(x).
\]
Из $P_x(U)=Q_x(U)-ax^4U$ и $a^2+b=0$ получаем тождество
\[
 H_N=H(x,p_N+x^Nz,y)
     =-ax^4y+Q_x(ax^{N+1}z+axp_N+y).
\]
Введём
\[
 \begin{aligned}
 G_N&=x^{N+1}z+by+x\bigl(p_N(x,y)-p_N(x,0)\bigr),\\
 c_N&=axp_N(x,0).
 \end{aligned}
\]
Тогда
\begin{equation}\label{eq:family-reduced}
 H_N-Q_x(c_N)
 =-ax^4y+\bigl[Q_x(aG_N+c_N)-Q_x(c_N)\bigr].
\end{equation}
Это приведённая запись из леммы~\ref{lem:inner-wild}:
$p_1=x^{N+1}$, $p_2=-ax^4$,
$Q_1(y)=by+x(p_N(x,y)-p_N(x,0))$.
Линейный коэффициент $Q_1$ равен $b$ по модулю $x$, поэтому
условие взаимной простоты выполнено. Внешний полином имеет степень
два, а $\deg_yQ_1\ge2$, поскольку $p_N\equiv2-y^2\pmod x$.

По модулю $x^{N+1}$ линейный коэффициент $Q_1$ обратим, а все
коэффициенты степеней $y^j$, $j\ge2$, нильпотентны.
Критерий Рассела--Сатайе показывает, что $G_N$ является
$R$-координатой. Но
\[
 [y^2]Q_1=-x\pmod{x^2},\qquad
 x^{N+1}\nmid[y^2]Q_1\quad(N\ge1).
\]
По критерию длины один $G_N$ дикая над $R$.
Лемма~\ref{lem:inner-wild} исключает ручность
$H_N-Q_x(c_N)$, а вычитание полинома от $x$ ручность не меняет.
Теорема Умирбаева--Ю даёт дикость $H_N$ в трёх переменных.
Поскольку $H_N$ входит в $\Phi_N$, последняя также дикая.
\end{proof}

Это доказательство относится ко всем $N$, а не выводится из проверки
конечного числа индексов. При $N=1$ его внутренний полином равен
$G_1=x^2z+by-xy^2$. Компьютерные проверки нормировки и коэффициентов
служат проверкой выкладок, но не заменяют лемму~\ref{lem:inner-wild}.

\section{Стандартная кремонова подгруппа}\label{sec:standard-subgroup}

Зафиксируем
\[
 \mathcal C_3(K)=\langle\operatorname{PGL}_4(K),\sigma_3\rangle,
 \qquad
 \sigma_3[X_0:X_1:X_2:X_3]
 =[X_1X_2X_3:X_0X_2X_3:X_0X_1X_3:X_0X_1X_2].
\]
В этом разделе произведения карт действуют справа налево.
По Блану--Хедену, $\Tame_3(K)\subset\mathcal C_3(K)$.
Их мономиальный критерий в размерности три утверждает, что
мономиальная бирациональная карта с матрицей показателей $M$ входит
в $\mathcal C_3(K)$ тогда и только тогда, когда суммы всех столбцов
$M$ нечётны \cite[теорема~1.1 и предложение~6.5]{BH2015}.

\begin{proposition}\label{prop:cremona-membership}
Все построенные выпрямители $\Phi_T,\Phi_H,\Phi_N$ принадлежат
$\mathcal C_3(K)$.
\end{proposition}
\begin{proof}
Рассмотрим подгруппу $\mathcal L\subset\Cr_3(K)$,
фиксирующую $x$ и порождённую картами
\[
 \begin{aligned}
 E_y(q)&=(x,y+q(x,z),z),& E_z(q)&=(x,y,z+q(x,y)),\\
 S&=(x,z,y),&B_r&=(x,x^ry,x^{-r}z),
 \end{aligned}
\]
где $q\in K[x,x^{-1},T]$, $r\in\Z$, и постоянными
диагональными картами $(x,cy,dz)$, $c,d\in K^\times$.

Все генераторы входят в $\mathcal C_3(K)$.
Для $B_r$ суммы столбцов матрицы показателей равны $(1,1,1)$.
Для $D=(x,x^{2s}y,z)$ они равны $(1+2s,1,1)$.
Выбрав $s\ge0$ так, чтобы $x^{2s}q$ был полиномом, получим
\[
 E_y(q)=D^{-1}E_y(x^{2s}q)D.
\]
Средняя карта ручная. Карты $E_z(q)$ получаются сопряжением с $S$;
остальные постоянные генераторы линейны.

В переменных $(u,v)$, при неизменном $x$, базовая карта раскладывается как
\[
 (u,v)\longmapsto\left(P_x(u)+mv,
 \frac{u-F_x(P_x(u)+mv)}m\right)
 =T_FD_mST_{P/m},
\]
где
\[
 \begin{aligned}
 T_{P/m}(u,v)&=(u,v+P_x(u)/m),&D_m(u,v)&=(mu,v/m),\\
 T_F(u,v)&=(u,v-F_x(u)/m),&S(u,v)&=(v,u).
 \end{aligned}
\]
При $m=bx^5$ каждый множитель принадлежит $\mathcal L$.
Начальная замена $(y,z)\mapsto(axy+z,-y)$ также является произведением
генераторов $\mathcal L$. Следовательно, $A_0\in\mathcal L$.

Для переноса положим $D_N=(x,y,x^Nz)$ и используем исправляющий
полином $\varphi_N(x,T)$ из конструкции полного компаньона. Тогда
\[
 \beta_N=SE_z(p_N)D_N,\qquad
 A_N=D_N^{-1}\bigl(E_z(-\varphi_N)A_0SE_z(p_N)\bigr)D_N.
\]
Карта $D_N$ нормализует $\mathcal L$, поскольку
\[
 \begin{aligned}
 D_N^{-1}E_y(q)D_N&=E_y(q(x,x^NT)),\\
 D_N^{-1}E_z(q)D_N&=E_z(x^{-N}q),\\
 D_N^{-1}SD_N&=B_NS,
 \end{aligned}
\]
а с $B_r$ и постоянными диагональными картами она коммутирует.
Те же формулы с $-N$ дают обратное включение.
Поэтому $A_N\in\mathcal L\subset\mathcal C_3(K)$.
Внешние карты, завершающие выпрямления, ручные; утверждение следует.
\end{proof}

\begin{remark}
При нечётном $N$ сама карта $D_N$ не проходит мономиальный критерий:
первая сумма столбца равна $1+N$. В доказательстве используется
\emph{нормализация} $\mathcal L$ картой $D_N$, а не принадлежность
$D_N$ группе $\mathcal C_3(K)$.
\end{remark}

\section{Стабильная ручность и границы вывода}\label{sec:stable-tameness}

\begin{proposition}\label{prop:stable-tameness}
Для каждого из построенных выпрямителей
\[
 \Phi_\ast\times\id_{\A^3}\in\Tame_6(K),
 \qquad \ast=T,H,N\ (N\ge1).
\]
\end{proposition}
\begin{proof}
Берсон--ван ден Эссен--Райт доказали, что подгруппа
$\Aut_K(\A^3)$, порождённая автоморфизмами, фиксирующими одну
координату, после добавления трёх переменных содержится в
$\Tame_6(K)$ в характеристике нуль
\cite[следствие~4.9]{BEW2012}.
Каждый $A_0,A_N$ фиксирует $x$.
Внешние множители наших выпрямителей ручные. Элементарные
автоморфизмы фиксируют две координаты, а аффинная группа порождена
элементарными линейными преобразованиями, масштабированиями одной
координаты и переносами одной координаты; каждый такой генератор фиксирует хотя бы одну
координату. Поэтому все $\Phi_\ast$ входят в указанную подгруппу.
\end{proof}

Оценка в три дополнительные переменные является достаточной;
минимальное число здесь не устанавливается. Объединяя доказательства,
получаем рабочий итог
\[
 \Phi_\ast\in
 \bigl(\mathcal C_3(K)\cap\Aut_K(\A^3)\bigr)
 \setminus\Tame_3(K),\qquad
 \Phi_\ast\times\id_{\A^3}\in\Tame_6(K).
\]
Это свойства \emph{выбранных карт}. Существование другого ручного
выпрямителя той же кривой ими не исключается.

Существование диких автоморфизмов в трёх переменных и диких координат
уже установлено в литературе. Принадлежность дикого автоморфизма Нагаты
группе $\mathcal C_3$ также доказана ранее
\cite[предложение~6.8]{BH2015}. Поэтому эти результаты не заявляются
как решение общей проблемы ручности, общей гипотезы стабильной ручности
или задачи о порождающих всей группы $\Cr_3$.
Содержание данного применения состоит в одновременной характеристике
явных выпрямителей конкретных кривых и в доказательстве для всех
индексов построенного семейства.

\section{Воспроизводимость и пределы вычислительного сертификата}
\label{sec:reproducibility}

\paragraph{Точная арифметика.}
Новый код использует Python и SymPy с рациональными коэффициентами
и явным соотношением $a^2-a+1=0$. Числа с плавающей точкой не
используются. Комментарии, сообщения об ошибках и описания функций
написаны по-английски. Проверки основаны на явных исключениях и
остаются включёнными при запуске Python с флагом \texttt{-O}.

\paragraph{Что проверяется вычислением.}
\begin{enumerate}
\item Обе конгруэнции для $P_x,F_x$, точные деления, компактная
базовая обратная карта и подстановки в трилистник.
\item Гармоническое тождество, перенесённая пара, точные деления
в прямом и обратном направлениях и восстановление кривой из оси.
\item Для явно указанного конечного диапазона индексов:
рекурсия $p_N$, полиномиальность $C_N$, обратимость центра по модулю
полного $x^N$, построение поправки, сертификаты полиномиальности
обеих карт и точные подстановки в параметризации.
\item Нормировка внутренней координаты в доказательстве дикости,
порядок её квадратичного коэффициента по $x$ и мономиальная чётность,
используемая в кремоновом доказательстве.
\item Отрицательные контрольные примеры: испорченная конгруэнция
и числитель, не допускающий требуемого точного деления, должны быть
обнаружены, включая режим \texttt{-O}.
\end{enumerate}

\paragraph{Граница вычислительного сертификата.}
Для больших карт код использует факторизованные сертификаты:
полиномиальность обеих сторон и тождества после локализации. Это
эквивалентно доказательству глобальной обратимости в области целостности;
полное раскрытие всех шести композиций в один огромный полином
не требуется и не заявляется для каждого $N$.
Утверждение для всех $N$ следует из индукции и доказательств в тексте,
а не из конечного диапазона тестов. Общие теоремы о дикости,
стабильной ручности и мономиальных преобразованиях используются
как опубликованные математические результаты. Прилагаемые программы
не являются их формальной верификацией в системе доказательств.
Явное слово в элементарных автоморфизмах шести переменных здесь
не вычислялось: существование такого слова следует из \cite{BEW2012}.

\paragraph{Запуск нового приложения.}
Из каталога \texttt{appendices/2026-10-28-cremona-program}:
\begin{verbatim}
python -m pip install -r requirements.txt
python computations/run_all.py
\end{verbatim}
Команда запускает обычный и оптимизированный режимы; точный диапазон,
версии программ, результаты и времена записываются в \texttt{results/}.
В поставляемом запуске проверены $N=1,2,3$: выполнено по 90 проверок
в каждом режиме, и их математические результаты совпали.
Описание отдельных проверок дано в \texttt{computations/README.md}.
Загружать данные из сети во время проверки не требуется.

\paragraph{Повторная проверка исходного препринта.}
Из корня полного репозитория выполнены
\begin{verbatim}
python -u verify.py
python -u verify_printed_formulas.py
\end{verbatim}
Оба запуска завершились успешно. Второй читает реальные правые части
формул исходного \LaTeX{} и проверяет все двенадцать компонент
прямых и обратных выпрямителей. Журналы находятся в
\texttt{results/upstream\_compact.log} и
\texttt{results/upstream\_printed.log} нового приложения.
Остальные старые журналы включены в полный архив как исторические
материалы; их наличие не означает, что каждый старый поиск, GAP-скрипт
или Macaulay2-скрипт был заново выполнен при подготовке Appendix 2.

\paragraph{Сборка и аудит файлов.}
Для \LaTeX{} используется \texttt{latex/main.tex}, компилятор XeLaTeX
и два прохода. Архив для Overleaf содержит все разделы, библиографию
и необходимые текстовые шрифты с лицензией. Преамбула сохраняет
формат исходного препринта: \texttt{article}, 11\,pt, A4, поля 24\,mm,
математические окружения, колонтитулы и кликабельные ссылки.
Файлы \texttt{PROVENANCE.md}, \texttt{audit/claim\_status.md} и
\texttt{SHA256SUMS.txt} описывают исходную версию, границы проверки
и контрольные суммы. Исходный основной текст и Appendix 1 не
переписываются. Это отдельная исследовательская ветка для дальнейших
исправлений и совместной работы.

\begin{thebibliography}{99}
\small\setlength{\itemsep}{6pt}
\bibitem{AM1975}
S.~S.~Abhyankar, T.~T.~Moh. \emph{Embeddings of the line in the plane}.
J. Reine Angew. Math. \textbf{276} (1975), 148--166.
\href{https://doi.org/10.1515/crll.1975.276.148}{doi:10.1515/crll.1975.276.148}.
\bibitem{Suzuki1974}
M.~Suzuki. \emph{Propriétés topologiques des polynômes de deux variables
complexes, et automorphismes algébriques de l'espace $\C^2$}.
J. Math. Soc. Japan \textbf{26} (1974), 241--257.
\href{https://doi.org/10.2969/jmsj/02620241}{doi:10.2969/jmsj/02620241}.
\bibitem{Shastri1992}
A.~R.~Shastri. \emph{Polynomial representations of knots}.
Tohoku Math. J. (2) \textbf{44} (1992), 11--17.
\href{https://doi.org/10.2748/tmj/1178227371}{doi:10.2748/tmj/1178227371};
\href{https://www.jstage.jst.go.jp/article/tmj1949/44/1/44_1_11/_pdf}{открытый PDF}.
\S1; \S3, Example~1.
\bibitem{BD1993}
S.~M.~Bhatwadekar, A.~K.~Dutta. \emph{On residual variables and stably
polynomial algebras}. Comm. Algebra \textbf{21} (1993), 635--645.
\href{https://doi.org/10.1080/00927879308824585}{doi:10.1080/00927879308824585}.
Theorem~3.2; связь с существованием компаньона.
\bibitem{EV2001}
E.~Edo, S.~Vénéreau. \emph{Length 2 variables of $A[x,y]$ and transfer}.
Ann. Polon. Math. \textbf{76} (2001), 67--76.
\href{https://doi.org/10.4064/ap76-1-6}{doi:10.4064/ap76-1-6};
\href{https://www.impan.pl/shop/publication/transaction/download/product/85204}{открытый PDF}.
Theorems~1--4, Proposition~2, \S5.
\bibitem{VdER2002}
A.~van den Essen, P.~van Rossum. \emph{Triangular derivations related
to problems on affine $n$-space}. Proc. Amer. Math. Soc.
\textbf{130} (2002), 1311--1322.
\href{https://doi.org/10.1090/S0002-9939-01-06212-8}{doi:10.1090/S0002-9939-01-06212-8}.
Theorem~4.1, Corollary~4.2, Conjectures~5.1--5.3.
\bibitem{UY2004}
U.~U.~Umirbaev, J.-T.~Yu. \emph{The strong Nagata conjecture}.
Proc. Natl. Acad. Sci. USA \textbf{101} (2004), 4352--4355.
\href{https://doi.org/10.1073/pnas.0308157101}{doi:10.1073/pnas.0308157101}.
Main Theorem; Theorems~5--6.
\bibitem{MP2009}
M.~Mella, E.~Polastri. \emph{Equivalent birational embeddings}.
Bull. London Math. Soc. \textbf{41} (2009), 89--93.
\href{https://doi.org/10.1112/blms/bdn107}{doi:10.1112/blms/bdn107};
\href{https://arxiv.org/abs/0906.4858}{arXiv:0906.4858}.
Theorem~1 and Corollary~1.
\bibitem{BEW2012}
J.~Berson, A.~van den Essen, D.~Wright.
\emph{Stable tameness of two-dimensional polynomial automorphisms over a regular ring}.
Adv. Math. \textbf{230} (2012), 2176--2197.
\href{https://doi.org/10.1016/j.aim.2012.04.017}{doi:10.1016/j.aim.2012.04.017};
\href{https://arxiv.org/abs/0707.3151}{arXiv:0707.3151}.
Corollary~4.9.
\bibitem{Edo2013}
E.~Edo. \emph{Coordinates of $R[x,y]$: Constructions and Classifications}.
Comm. Algebra \textbf{41} (2013), 4694--4710.
\href{https://doi.org/10.1080/00927872.2012.700981}{doi:10.1080/00927872.2012.700981};
\href{https://arxiv.org/abs/1206.5876}{arXiv:1206.5876}.
Proposition~2, Theorem~8(1) and proof of Theorem~8(3).
\bibitem{BH2015}
J.~Blanc, I.~Hedén. \emph{The group of Cremona transformations generated
by linear maps and the standard involution}.
Ann. Inst. Fourier \textbf{65} (2015), 2641--2680.
\href{https://doi.org/10.5802/aif.2999}{doi:10.5802/aif.2999};
\href{https://www.numdam.org/item/10.5802/aif.2999.pdf}{открытый PDF}.
Theorem~1.1; Propositions~6.5, 6.8.
\bibitem{Lewis2024}
D.~Lewis. \emph{A unified approach to embeddings of a line in 3-space}.
Comm. Algebra \textbf{52} (2024), 189--200.
\href{https://doi.org/10.1080/00927872.2023.2239344}{doi:10.1080/00927872.2023.2239344};
\href{https://arxiv.org/abs/2203.08178}{arXiv:2203.08178} (препринт 2022).
Lemma~1, Theorems~2, 12; введение о предшествующих результатах.
\bibitem{OdaWithdrawn}
S.~Oda. \emph{Note on the candidate counter-example in the cancellation
problem for affine spaces posed by Arno Van den Essen}.
\href{https://arxiv.org/abs/1201.4198}{arXiv:1201.4198}.
Отозван автором 9 мая 2020 года; не используется как доказательство.
\bibitem{Repository2026}
I.~Dorofeev. \emph{Polynomial rectification of the trefoil and a harmonic
figure-eight}. Исследовательский репозиторий, 2026.
\href{https://github.com/qft98romanov-hash/abhyankar-moh-trefoil/tree/b954f3609e900761fddc9c8ff03d8301e9604cab}{Зафиксированная исходная версия};
\href{https://github.com/qft98romanov-hash/abhyankar-moh-trefoil/tree/b954f3609e900761fddc9c8ff03d8301e9604cab/appendices/2026-10-26-harmonic-lnd}{Appendix 1}.
Материалы проекта, а не внешняя рецензированная публикация.
\end{thebibliography}

\end{document}
